\documentclass[11pt]{article}
% Standalone preamble for web-only supplementary notes.
\usepackage[T1]{fontenc}
\usepackage{amsmath,amssymb,amsthm}
\usepackage{hyperref}
\usepackage{cleveref}
\DeclareMathOperator{\Map}{Map}
\DeclareMathOperator{\Fun}{Fun}
\newcommand{\iso}{\xrightarrow{\sim}}
\newtheorem{lemma}{Supplementary Lemma}
\newtheorem{corollary}[lemma]{Supplementary Corollary}
\crefname{lemma}{Supplementary Lemma}{Supplementary Lemmas}
\Crefname{lemma}{Supplementary Lemma}{Supplementary Lemmas}
\crefname{corollary}{Supplementary Corollary}{Supplementary Corollaries}
\Crefname{corollary}{Supplementary Corollary}{Supplementary Corollaries}

\begin{document}
\section*{Arrows with specified endpoints}

The hom and slice constructions in \href{morphisms-and-diagrams.html}{Section 2.1}
use a common pullback of the endpoint functor. This note explains how the
endpoint identifications are recovered from its matching, and how a family
of arrows gives a functor into that pullback.

Write $\operatorname{Ar}(C)=\Fun([1],C)$ and let $e_0,e_1$ be the two
evaluation functors. For $u,v\colon B\to C$, put
\[
 F(u,v)=\operatorname{Ar}(C)\times_{C\times C}B,
\]
where the maps to $C\times C$ are $(e_0,e_1)$ and $(u,v)$.
Denote its projections by $a$ and $b$, and its specified matching by
$\mu\colon(e_0,e_1)a=(u,v)b$. All categories here are absolute;
a parameter category $\Gamma$ describes a functor into these fixed categories.

\begin{lemma}
\label[lemma]{lem:Supplementary_Endpoint_Frames}
The matching $\mu$ determines specified identifications
\[
 p_0\colon e_0a=ub,\qquad p_1\colon e_1a=vb.
\]
They are obtained by projecting $\mu$ and composing with the projection
comparisons for the pairings on its two sides.
\end{lemma}
\begin{proof}[Proof idea]
For the first coordinate, let $L_0\colon\pi_1((e_0,e_1)a)=e_0a$
and $R_0\colon\pi_1((u,v)b)=ub$ be the product comparisons, including
reassociation. Then
\[
 p_0=R_0\cdot(\pi_1\star\mu)\cdot L_0^{-1}.
\]
The second coordinate has the same formula using $\pi_2$.
Here $\star$ denotes whiskering and vertical composition applies the
rightmost identification first. These formulas give the source and target
frames of the universal arrow family.
\end{proof}

\begin{lemma}
\label[lemma]{lem:Supplementary_Endpoint_Lift}
Suppose $y\colon\Gamma\to B$ and $A\colon\Gamma\to\operatorname{Ar}(C)$
are equipped with identifications $s\colon e_0A=uy$ and
$t\colon e_1A=vy$. There is a functor $\ell\colon\Gamma\to F(u,v)$,
with comparisons $a\ell=A$ and $b\ell=y$, compatible with the specified
endpoint identifications $s,t$.
\end{lemma}
\begin{proof}[Proof idea]
Write $\kappa(f,g,r)\colon(f,g)r=(fr,gr)$ for the pairing and
precomposition comparison. The cone matching determined by $s,t$ is
\[
 (e_0,e_1)A
 \xrightarrow{\kappa(e_0,e_1,A)}(e_0A,e_1A)
 \xrightarrow{(s,t)}(uy,vy)
 \xrightarrow{\kappa(u,v,y)^{-1}}(u,v)y.
\]
Apply the pullback lifting operation to this cone. Its computation law
compares the entire induced cone with the displayed one. In particular,
it supplies both projection comparisons and their compatibility with
the matching. Project that compatibility to recover the two endpoint equations.
\end{proof}

For absolute objects $x,y$ of $C$, take $B=\star$, $u=x$, and $v=y$.
The resulting endpoint fiber is the hom category $\operatorname{Hom}_C(x,y)$.
Taking $B=C$, one of $u,v$ to be the identity, and the other to be
constant at a chosen object gives the corresponding slice construction. In each case, an input includes the identifications
with the specified endpoints. The full lifting comparison is what permits
subsequent calculations to use those same endpoints.

The normalization of the terminal parameter in the hom case is explained
further in \href{supplementary-hom-expressions-and-restriction.html}{Hom expressions and restriction}.
\end{document}
